% ============================================================================
% 07_humans.tex — Ordinal Agreement and the Limits of Cardinal Calibration
% (sec:humans). Rewritten in the merge pass: Part 1 = rung-by-rung ordinal
% anchoring (updated to the four-model E3 cells and the canonical direction
% convention); Part 2 = the cardinal-calibration question (tab:cardinal),
% built from results/cardinal/{cardinal_grid.csv,cardinal_summary.md}.
% Human cells are moment-matched reconstructions (sec:reconstruction, app:human).
% ============================================================================

\section{Ordinal Agreement and the Limits of Cardinal Calibration}\label{sec:humans}

The ranking of Section~\ref{sec:ranking} is estimated on populations of LLM agents. This section asks, in two steps, what that ranking is worth to a designer whose participants are human---and, symmetrically, what five decades of human experimental evidence contribute to the credibility of the ranking. The first step is ordinal: on every rung of the ranking where a human benchmark exists, the sign of the treatment effect agrees, even though the two populations make opposite baseline mistakes---a fact we argue converts the validity map's most awkward cell into the strongest available evidence that the ranking measures properties of \emph{mechanisms} rather than of any particular reasoner. The second step is cardinal, and its answer is negative in an instructive way: searching a grid of the tuning knobs available to a practitioner---sampling temperature, risk personas, framing manipulations, rule explanations---we find that no single setting produces cardinal agreement with human play across mechanisms. Level agreement is achievable mechanism by mechanism; direction agreement is achievable with exactly one knob that breaks the level agreement elsewhere; and the best tuning for one model family overshoots badly on another. Together the two steps justify the discipline imposed since Section~\ref{sec:reconstruction}: the testbed is an instrument for signs, orderings, and comparative statics, and the boundary of that instrument is now measured rather than assumed.

\subsection{Rung-by-Rung Ordinal Agreement}\label{sec:humans-rungs}

We walk down the ladder of Table~\ref{tab:ranking}, rung by rung, comparing each lever's effect on LLM play with its effect on human subjects where a human experiment exists. The human-anchored rungs are marked as diamonds in Figure~\ref{fig:ranking}; human quantities below are computed from the moment-matched reconstructions of Section~\ref{sec:reconstruction} and anchored to the original published findings.

\paragraph{Extensive form (Family A).} Clock implementations help humans. \citet{li2017obviously} showed experimentally that ascending clock auctions, which are obviously strategy-proof, produce markedly more dominant-strategy play than the strategically equivalent sealed-bid format, and \citet{breitmoser2022obviousness} replicate and extend this pattern; in our reconstruction, human SMAD falls from $9.31\%$ under the sealed-bid second-price format to $3.54\%$ under AC and $5.83\%$ under AC-B, for a preservation index of $+0.62$ $[+0.39,+0.79]$ (AC) and $+0.37$ $[-0.12,+0.72]$ (AC-B). The same lever helps every model we test, and helps it more: in the harmonized four-model grid, the closed clock cuts SMAD to $1.4$--$5.2\%$ from sealed-bid baselines of $9.5$--$20.7\%$ ($p<0.001$ for all four models; Cohen's $d$ up to $1.00$; Section~\ref{sec:ranking-extensive}). The second Family-A comparison comes from the matching robustness domain (Appendix~\ref{app:da}): the iterative, obviously strategy-proof implementation of DA eliminated every observable misreport for three of the four models,\footnote{The zero is exact for the pick-based query protocol in three of four models (rule-of-three $95\%$ upper bounds of $0.85$--$0.92\%$ on the per-decision error rate); Gemma commits five strict pick errors ($1.6\%$), and an alternative yes/no-question implementation of the same extensive form exposes model-specific failures, most notably Claude's $24.1\%$ rate of false rejections. The rung's zero is a statement about a protocol, not about cognition; Appendix~\ref{app:da} gives the full accounting.} and, to our knowledge, no experiment has tested an OSP implementation of DA with human subjects. That rung is issued as a prediction rather than an anchor, and we return to unanchored rungs below.

\paragraph{Clock-framing (B3).} \citet{breitmoser2022obviousness} show that \emph{describing} the sealed-bid second-price auction as an automatic exit threshold in an ascending clock---changing words, not the mechanism---substantially increases truthful bidding among human subjects. The same one-paragraph description now has four-model support: it reduces SMAD in every family---Claude $9.5\%\to7.2\%$, GPT-4o $11.4\%\to6.3\%$, and Gemma $20.7\%\to13.8\%$ (each $p<0.001$), with Gemini's improvement smaller and not significant ($11.7\%\to9.5\%$, $p=0.097$)---for a pooled preservation index of $+0.29$ $[+0.17,+0.45]$ (Section~\ref{sec:ranking-descriptions}). The sign agrees with the human finding in all four families; the magnitude is uniformly positive but well below the true clock's.

\paragraph{Menu restatement (B1).} In the laboratory experiments of \citet{gonczarowski2023strategyproofness}, a menu description that compresses the dominant-strategy logic of the second-price auction does not meaningfully change human bidding (reconstructed $\rho=-0.07$ $[-0.12,-0.02]$). Across our four model families the restatement has no consistent effect either, but the agreement is an agreement \emph{on average} rather than cell by cell: the effect is null for Gemma ($p=0.94$), moves Claude slightly away from truthful play ($+\$0.88$, $p=0.042$) and GPT-4o more clearly away ($-\$1.32$, $p=0.003$), yet moves Gemini \emph{toward} truthful play ($+\$1.75$, $p<0.001$). Pooling, the effects are small and sign-mixed, and a two-sided sign test cannot reject a zero median ($p=0.73$; Section~\ref{sec:ranking-descriptions})---in sharp contrast to the uniformly signed effects of safety-exposing text. We therefore state this rung's anchor precisely: LLM play is consistent with the human null \emph{on average}, with model-level heterogeneity that the single human point estimate cannot adjudicate. The rung retains its placebo value in one direction---the testbed does not manufacture a reliable gain from restating the mechanics---but it is not the clean, uniformly inert tier the human evidence might suggest.

\paragraph{Baseline.} Both populations deviate under the direct mechanisms. Human subjects systematically overbid in sealed-bid second-price auctions \citep{kagel1993independent, kagel1995individual} and submit dominated rank-order lists in deployed matching mechanisms \citep{rees2018suboptimal}; all four LLM families underbid in SPSB, and all four misreport in the direct-DA robustness domain (Section~\ref{sec:validation}, Section~\ref{sec:ranking}, Appendix~\ref{app:da}).

In sum: the clock helps both populations, the clock-framing description helps both, the menu restatement moves neither on average, and the baseline direct mechanisms fail both. Every rung with a human benchmark agrees in sign, with the menu agreement holding at the pooled level over sign-mixed model cells.

\subsection{Opposite Failures, Common Levers: Substrate Independence}\label{sec:humans-reversal}

The validity map's red cell is the direction reversal: humans overbid in SPSB-IPV---$67.2\%$ of all bids exceed value, and $92.2\%$ of the bids that deviate from value deviate upward---while GPT-4o's deviations are overwhelmingly ($81.6\%$) underbids; the same reversal appears in the APV sealed-bid cell ($64.7\%$ human overbidding against $92.4\%$ GPT-4o underbidding).\footnote{Direction shares depend on the convention. Throughout this section we use the repo-canonical convention---the share of bids strictly above value among \emph{all} bids---under which the human SPSB overbid share is $67.2\%$; conditioning on bids that deviate from value gives $92.2\%$. The directional decomposition of Section~\ref{sec:validation}, which removes a $\pm2\%$ band around the benchmark before classifying, yields slightly different deviator-conditional shares; no conclusion in this paper turns on the choice.} Read as a test of the proxy hypothesis---``LLMs reproduce human behavior''---this cell is a failure, and we declared it as such in Section~\ref{sec:validation}. Read against the ranking, it becomes the paper's sharpest identification result. Two populations fail in \emph{opposite} directions, yet they are corrected by the \emph{same} levers, in the same order: the clock repairs human overbidding and LLM underbidding alike, the clock-framing description does the same, and the menu restatement repairs neither. If the deviations were driven by population-specific psychology---a joy of winning or spite pushing humans above value, a trained conservatism pulling LLMs below it---there would be no reason for a shared correction structure. The parsimonious reading is that the binding constraint lives in the strategic problem itself: in the difficulty of verifying, from a direct mechanism's description, that truth-telling is safe \citep{li2017obviously, oprea2024decisions}. What fills the gap left by that missing verification is population-specific---competitive overbidding in humans, conservative shading in LLMs---but the gap, and the levers that close it, are shared. This is the sense in which the lever ordering is substrate-independent, echoing the hypothesis structure of Section~\ref{sec:framework}.

The reversal also rules out a deflationary explanation of the rung-by-rung agreement. If LLMs merely reproduced human experimental regularities memorized from training corpora, agreement on the levers would be uninformative---but a population imitating the published human data would \emph{overbid} in second-price auctions, and ours does not. The agreement on which levers work is therefore not inherited from the human literature; it is re-derived by a different reasoner failing in a different way.

\paragraph{An out-of-sample prediction.} One tier of the ranking has no human anchor precisely where the ranking is most useful. Safety-exposing descriptions (B2) are the cheapest lever in the designer's toolkit and, on LLM populations, among the most powerful; and the robustness domain sharpens the prediction into a contrast, since there the \emph{same} menu format splits by content---a menu statement that includes the invariance property behaves like a safety description, while a menu that merely restates the outcome computation backfires (Appendix~\ref{app:da}). The closest human evidence is instructive but incomplete: \citet{gonczarowski2024describing} find that a menu-based explanation of DA significantly improved participants' measured \emph{understanding} of strategy-proofness, while average \emph{behavioral} effects were modest. Our decomposition predicts where the behavioral gain is: in descriptions that expose the safety invariant (``your bid determines whether you win, not what you pay''; ``rejections only redirect, never eliminate''), not in descriptions that restate the outcome mapping. Concretely, we predict that in matched human experiments, a Payoff Safety description will shift human SPSB bid distributions toward value where the menu framing of \citet{gonczarowski2023strategyproofness} did not, and that invariance-stating explanations of DA will move behavior where mechanics-restating explanations move only comprehension. The prediction is cheap to test and falsifiable in both directions: the synthetic version of each treatment cell cost roughly \$10 to run, and an incentivized human counterpart is a single-session design at roughly the \$2{,}000 scale of existing description experiments (Table~\ref{tab:cost}).

\subsection{From Ordinal to Cardinal: Is There a Consistent Tuning?}\label{sec:humans-cardinal}

Everything so far is ordinal, and deliberately so: where levels can be compared, they diverge. LLMs \emph{over-comply} in exactly the formats where the levers work best---in the clock implementations, model SMAD of $0.5$--$0.8\%$ sits an order of magnitude below the human $3.5$--$5.8\%$---while in third-price auctions they sit far above any human benchmark. A natural response, and an increasingly common proposal in the literature on LLMs as experimental proxies, is to \emph{tune}: adjust temperature, personas, or framings until the synthetic population matches the human one cardinally, then read levels off the tuned simulator. This subsection asks whether any such tuning exists in our environment, using every knob for which we have data: sampling temperature ($T\in\{0.1,0.5,1.0\}$; GPT-4o, eight mechanisms), risk personas (averse/neutral/seeking; second price for all four models, first and third price for GPT-4o), prospect-theoretic framings (loss/gain/mixed/endowment/WTA--WTP), and rule explanation on/off---a grid of 95 mechanism$\times$knob cells. Distance to the human benchmarks is measured on three axes: \emph{level} (the gap in SMAD), \emph{direction} (the gap in the share of bids above value), and \emph{shape} (Wasserstein-1 distance between the scaled deviation distributions and the moment-matched human reconstructions). Table~\ref{tab:cardinal} reports the core sealed-bid grid.

\begin{table}[t]
\centering
\small
\caption{Cardinal calibration under tuning knobs (GPT-4o, sealed-bid IPV). Each cell reports SMAD\,\% and the share of bids above value; human targets in the top row. SMAD $=100\cdot\mathbb{E}|b-b^{*}(v)|/24.5$ with each environment's own equilibrium $b^{*}$ (FPSB $N{=}3$: $2v/3$; SPSB: $v$; TPSB $N{=}3$: $2v$). $\dagger$: differs from its own baseline at $p<0.05$ (Welch tests for SMAD, two-proportion tests for shares); baselines are the pooled axis baselines (personas), the loss-aversion baseline (frames), and the $T{=}0.5$ anchor (temperature, explanation). The human third-price target is the \citet{kagel1993independent} $n{=}5$ environment, the closest available benchmark for the $N{=}3$ knob cells.}
\label{tab:cardinal}
\begin{tabular}{lcccccc}
\toprule
& \multicolumn{2}{c}{First price} & \multicolumn{2}{c}{Second price} & \multicolumn{2}{c}{Third price ($N{=}3$)} \\
Setting & SMAD & over\,\% & SMAD & over\,\% & SMAD & over\,\% \\
\midrule
\emph{Human target} & \emph{24.8} & \emph{0.4} & \emph{5.6} & \emph{67.2} & \emph{7.7}$^{\,n=5}$ & \emph{89.8} \\
\midrule
Untuned anchor ($T{=}0.5$)   & 28.1 & 0.0  & 8.9  & 16.7 & 108.3 & 19.4 \\
Axis-pooled baseline     & 23.4 & 0.0  & 11.4 & 0.4  & 120.2 & 0.2  \\
Temperature 0.1          & 28.6 & 0.0  & 15.1$^\dagger$ & 10.2$^\dagger$ & 112.5 & 1.3$^\dagger$ \\
Temperature 1.0          & 25.5 & 0.7  & 13.2$^\dagger$ & 9.8$^\dagger$  & 123.7$^\dagger$ & 2.7$^\dagger$ \\
Risk-averse persona      & 15.0$^\dagger$ & 0.0 & 18.3$^\dagger$ & 0.0 & 129.2 & 0.0 \\
Risk-neutral persona     & 20.9 & 0.0 & 11.0 & 2.7$^\dagger$ & 121.3 & 6.7$^\dagger$ \\
Risk-seeking persona     & 28.0$^\dagger$ & 18.0$^\dagger$ & 10.5 & 45.3$^\dagger$ & 105.8$^\dagger$ & 56.5$^\dagger$ \\
WTA--WTP frame           & 23.0 & 0.0  & 7.2$^\dagger$  & 0.0  & 114.0 & 0.0 \\
Rule explanation off     & 28.1 & 0.0  & 8.7  & 16.0 & 111.3 & 12.6$^\dagger$ \\
\bottomrule
\end{tabular}
\end{table}

\paragraph{Level: the best knob is mechanism-specific.} On the level axis alone, cardinal agreement is often achievable---but with a different knob in every format. In FPSB, temperature $1.0$ closes the gap almost exactly (SMAD $25.50$ $[23.49,27.69]$ against a human $24.76$; untuned anchor $28.11$). In SPSB, the WTA--WTP frame is best ($7.24$ $[6.02,8.56]$ against $5.65$; $-3.57$ SMAD relative to its own baseline, Welch $p=0.005$), while both temperature deviations from the anchor make play \emph{worse}. In the affiliated-values sealed-bid format the untuned anchor is already the best available cell ($10.95$ against $9.31$). The clock formats invert the problem: LLMs are \emph{too good}---SMAD $0.83$ (AC) and $0.48$ (AC-B) against human targets of $3.54$ and $5.83$---so the only knob that moves them toward humans is raising temperature to $1.0$, which works purely by injecting noise into near-perfect play.\footnote{The high-temperature closed-clock cell is small ($8$ non-winner exits); we read the clock rows as suggestive. Two grid cells are unavailable (both clock formats at $T{=}0.1$), and the third-price persona and frame cells exist only at $N{=}3$.} And third-price auctions are not tunable at all: the best third-price cell anywhere in the grid has SMAD $44.1$ (temperature $0.1$, $N{=}5$) against a human $7.66$, and at $N{=}3$---where the equilibrium bid is $2v$---every persona and frame cell sits at SMAD $106$--$131$, because the models bid approximately their value under every setting.

\paragraph{Direction: one knob moves it, at a price.} The direction axis is where the human--LLM reversal of Section~\ref{sec:humans-reversal} lives, and exactly one knob addresses it. Under the risk-seeking persona, GPT-4o's share of above-value bids rises from $0.4\%$ to $45.3\%$ in SPSB and from $0.2\%$ to $56.5\%$ in third price (both $p<10^{-4}$), against human targets of $67.2\%$ and $89.8\%$---direction gaps fall from $66.8$ to $21.9$ percentage points and from $89.6$ to $33.3$. The same persona also improves the third-price \emph{level} ($-14.4$ SMAD, $p=0.011$) and the second-price \emph{shape} (Wasserstein-1 from $23.8$ to $14.8$ points). No other knob moves direction: under every temperature, frame, and explanation setting, the SPSB overbid share never exceeds $17\%$, and the risk-averse persona drives all overbid shares to zero. But the price is paid in first price, the one format where humans almost never bid above value: the risk-seeking persona induces dominated bids in $18.0\%$ of FPSB decisions (human: $0.4\%$; $p<10^{-4}$) and raises FPSB SMAD by $4.6$ points ($p=0.002$). In the direction rankings the same setting is first of fifteen in both SPSB and third price and \emph{last} of fifteen in FPSB---the sharpest possible statement of the trade-off. Matching human overbidding where humans overbid requires a disposition that violates the sharpest regularity humans display where they do not.

\paragraph{Tuning is model-specific.} The trade-off does not even transfer across model families. The same risk-seeking persona applied to SPSB moves Gemma from $20.7$ to $5.87$ SMAD---against the human $5.65$, the single best cardinal match in the grid---and Gemini from $11.7$ to $4.9$, but leaves GPT-4o essentially unchanged ($11.4$ to $10.5$) and explodes Claude from $9.5$ to $35.0$, with $93\%$ of bids above value. A persona calibrated on one model family, then deployed on another, can convert mild underbidding into catastrophic overbidding.

\paragraph{The remaining knobs are not dials.} Temperature is non-monotone where it matters: SPSB SMAD is $15.1$, $8.9$, and $13.2$ at $T=0.1$, $0.5$, and $1.0$ respectively (both deviations from the anchor significant at $p\le10^{-4}$), so it cannot be used to trade off fidelity smoothly; its only clean cardinal use is adding human-like noise to the too-perfect clock formats. Turning the rule explanation on or off has no detectable effect on any axis (level differences of at most $0.16$ SMAD in the first-, second-, and affiliated-value formats, all $p\ge0.53$).\footnote{We phrase this as ``no detectable effect'' rather than as an estimated effect of removing explanation: the manipulation cannot be verified from the archived prompts, and the contrast rests on the run directories' provenance.} The remaining prospect-theoretic frames mostly \emph{hurt} the second-price level relative to their own baseline.

\paragraph{Verdict.} There is no single consistent tuning. Formally: no setting in the grid is in the top half of both the level ranking and the direction ranking on all three core mechanisms; the level winner by worst-case rank (the WTA--WTP frame) and the direction winner (the risk-seeking persona) are disjoint, and the direction winner breaks the level axis in first price. Cardinal agreement is therefore achievable only mechanism-by-mechanism, and sometimes only model-by-model---a boundary that is structural rather than an accident of our grid, since the direction axis demands a disposition that the first-price format punishes. Two implications follow. First, the boundary bounds the proxy use-case: a designer who needs calibrated \emph{levels}---revenue predictions, welfare magnitudes---cannot obtain them from a single tuned synthetic population that spans mechanisms, which is why every cross-population claim in this paper is ordinal. Second, it clarifies what \emph{would} be needed: per-mechanism (and, our results add, per-model) calibration against a small sample of human data from the target environment, in the spirit of \citet{corrigan2025usefulness}---a local correction that our grid suggests cannot be amortized across formats, because the knob that fixes one format's distribution is the knob that breaks another's.
